\chapter{Integral Dependence and Valuations}\label{chapter5}

In classical algebraic geometry curves were frequently studied by projecting them onto a line and regarding the curve as a (ramified) covering of the line. This is quite analogous to the relationship between a number field and the rational field---or rather between their rings of integers---and the common algebraic feature is the notion of integral dependence. In this chapter we prove a number of results about integral dependence. In particular we prove the theorems of Cohen-Seidenberg (the ``going-up'' and ``going-down'' theorems) concerning prime ideals in an integral extension. In the exercises at the end we discuss the algebro-geometric situation and in particular the Normalization Lemma.

We also give a brief treatment of valuations.

\section*{Integral Dependence}
\addcontentsline{toc}{section}{Integral Dependence}
Let $B$ be a ring, $A$ a subring of $B$ (so that $1 \in A$). An element $x$ of $B$ is said to be \emph{integral over $A$} if $x$ is a root of a monic polynomial with coefficients in $A$, that is if $x$ satisfies an equation of the form
\begin{equation}\label{eq:integral-eq}
x^n + a_1 x^{n-1} + \cdots + a_n = 0 \tag{1}
\end{equation}
where the $a_i$ are elements of $A$. Clearly every element of $A$ is integral over $A$.

\setcounter{theorem}{-1}
\begin{examplenumbered}\label{ex:5.0}
$A = \Z$, $B = \Q$. If a rational number $x = r/s$ is integral over $\Z$, where $r, s$ have no common factor, we have from \eqref{eq:integral-eq}
\[
r^n + a_1 r^{n-1} s + \cdots + a_n s^n = 0
\]
the $a_i$ being rational integers. Hence $s$ divides $r^n$, hence $s = \pm 1$, hence $x \in \Z$.
\end{examplenumbered}

\begin{proposition}\label{prop:5.1}\index{integral element!characterization}
The following are equivalent:
\begin{enumerate}
    \item[i)] $x \in B$ is integral over $A$;
    \item[ii)] $A[x]$ is a finitely generated $A$-module;
    \item[iii)] $A[x]$ is contained in a subring $C$ of $B$ such that $C$ is a finitely generated $A$-module;
    \item[iv)] There exists a faithful $A[x]$-module $M$ which is finitely generated as an $A$-module.
\end{enumerate}
\end{proposition}

\begin{proof}
    i) $\implies$ ii). From \eqref{eq:integral-eq} we have
    \[
    x^{n+r} = -a_1 x^{n+r-1} - \cdots - a_n x^r
    \]
    for all $r \geq 0$; hence, by induction, all positive powers of $x$ lie in the $A$-module generated by $1, x, \ldots, x^{n-1}$. Hence $A[x]$ is generated (as an $A$-module) by $1, x, \ldots, x^{n-1}$.

    ii) $\implies$ iii). Take $C = A[x]$.

    iii) $\implies$ iv). Take $M = C$, which is a faithful $A[x]$-module (since $yC = 0 \implies y \cdot 1 = 0$).

    iv) $\implies$ i). This follows from \eqref{prop:module-endomorphism-equation}: take $\phi$ to be multiplication by $x$, and $\mathfrak{a} = A$ (we have $xM \subseteq M$ since $M$ is an $A[x]$-module); since $M$ is faithful, we have $x^n + a_1 x^{n-1} + \cdots + a_n = 0$ for suitable $a_i \in A$.
\end{proof}

\begin{corollary}\label{cor:5.2}
    Let $x_i$ ($1 \leq i \leq n$) be elements of $B$, each integral over $A$. Then the ring $A[x_1, \ldots, x_n]$ is a finitely-generated $A$-module.
\end{corollary}

\begin{proof}
    By induction on $n$. The case $n = 1$ is part of \eqref{prop:5.1}. Assume $n > 1$, let $A_r = A[x_1, \ldots, x_r]$; then by the inductive hypothesis $A_{n-1}$ is a finitely generated $A$-module. $A_n = A_{n-1}[x_n]$ is a finitely generated $A_{n-1}$-module (by the case $n = 1$, since $x_n$ is integral over $A_{n-1}$). Hence by \eqref{prop:finitely-generated-restriction}, $A_n$ is finitely generated as an $A$-module.
\end{proof}

\begin{corollary}\label{cor:5.3}\index{integral closure}
    The set $C$ of elements of $B$ which are integral over $A$ is a subring of $B$ containing $A$.
\end{corollary}

\begin{proof}
    If $x, y \in C$ then $A[x, y]$ is a finitely generated $A$-module by \eqref{cor:5.2}. Hence $x \pm y$ and $xy$ are integral over $A$, by (iii) of \eqref{prop:5.1}.
\end{proof}

The ring $C$ in \eqref{cor:5.3} is called the \emph{integral closure} of $A$ in $B$.\index{integral closure} If $C = A$, then $A$ is said to be \emph{integrally closed} in $B$.\index{integrally closed} If $C = B$, the ring $B$ is said to be \emph{integral} over $A$.

\begin{remark}
    Let $f \colon A \to B$ be a ring homomorphism, so that $B$ is an $A$-algebra. Then $f$ is said to be \emph{integral}, and $B$ is said to be an \emph{integral $A$-algebra}, if $B$ is integral over its subring $f(A)$. In this terminology, the above results show that
    \[
    \text{finite type} + \text{integral} = \text{finite}.
    \]
\end{remark}

\begin{corollary}\label{cor:5.4}
    If $A \subseteq B \subseteq C$ are rings and if $B$ is integral over $A$, and $C$ is integral over $B$, then $C$ is integral over $A$ (transitivity of integral dependence).
\end{corollary}

\begin{proof}
    Let $x \in C$, then we have an equation
    \[
    x^n + b_1 x^{n-1} + \cdots + b_n = 0 \qquad (b_i \in B).
    \]
    The ring $B' = A[b_1, \ldots, b_n]$ is a finitely generated $A$-module by \eqref{cor:5.2}, and $B'[x]$ is a finitely generated $B'$-module (since $x$ is integral over $B'$). Hence $B'[x]$ is a finitely generated $A$-module by \eqref{prop:finitely-generated-restriction} and therefore $x$ is integral over $A$ by (iii) of \eqref{prop:5.1}.
\end{proof}

\begin{corollary}\label{cor:5.5}
Let $A \subseteq B$ be rings and let $C$ be the integral closure of $A$ in $B$. Then $C$ is integrally closed in $B$.
\end{corollary}

\begin{proof}
    Let $x \in B$ be integral over $C$. By \eqref{cor:5.4}, $x$ is integral over $A$, hence $x \in C$.
\end{proof}

The next proposition shows that integral dependence is preserved on passing to quotients and to rings of fractions:

\begin{proposition}\label{prop:5.6}
Let $A \subseteq B$ be rings, $B$ integral over $A$.
\begin{enumerate}
    \item[i)] \label{item 1:prop:5.6} If $\mathfrak{b}$ is an ideal of $B$ and $\mathfrak{a} = \mathfrak{b}^c = A \cap \mathfrak{b}$, then $B/\mathfrak{b}$ is integral over $A/\mathfrak{a}$.
    \item[ii)] \label{item 2:prop:5.6} If $S$ is a multiplicatively closed subset of $A$, then $S^{-1}B$ is integral over $S^{-1}A$.
\end{enumerate}
\end{proposition}

\begin{proof}
i) If $x \in B$ we have, say, $x^n + a_1 x^{n-1} + \cdots + a_n = 0$, with $a_i \in A$. Reduce this equation mod~$\mathfrak{b}$.

ii) Let $x/s \in S^{-1}B$ ($x \in B$, $s \in S$). Then the equation above gives
\[
(x/s)^n + (a_1/s)(x/s)^{n-1} + \cdots + a_n/s^n = 0
\]
which shows that $x/s$ is integral over $S^{-1}A$.
\end{proof}

\section*{The Going-Up Theorem}
\addcontentsline{toc}{section}{The Going-Up Theorem}

\begin{proposition}\label{prop:5.7}\index{integral extension!field criterion}
Let $A \subseteq B$ be integral domains, $B$ integral over $A$. Then $B$ is a field if and only if $A$ is a field.
\end{proposition}

\begin{proof}
Suppose $A$ is a field; let $y \in B$, $y \neq 0$. Let
\[
y^n + a_1 y^{n-1} + \cdots + a_n = 0
\]
be an equation of integral dependence for $y$ of smallest possible degree. Since $B$ is an integral domain we have $a_n \neq 0$, hence $y^{-1} = -a_n^{-1}(y^{n-1} + a_1 y^{n-2} + \cdots + a_{n-1}) \in B$. Hence $B$ is a field.

Conversely, suppose $B$ is a field; let $x \in A$, $x \neq 0$. Then $x^{-1} \in B$, hence is integral over $A$, so that we have an equation
\[
x^{-m} + a'_1 x^{-m+1} + \cdots + a'_m = 0 \qquad (a'_i \in A).
\]
It follows that $x^{-1} = -(a'_1 + a'_2 x + \cdots + a'_m x^{m-1}) \in A$, hence $A$ is a field.
\end{proof}

\begin{corollary}\label{cor:5.8}
    Let $A \subseteq B$ be rings, $B$ integral over $A$; let $\mathfrak{q}$ be a prime ideal of $B$ and let $\mathfrak{p} = \mathfrak{q}^c = \mathfrak{q} \cap A$. Then $\mathfrak{q}$ is maximal if and only if $\mathfrak{p}$ is maximal.
\end{corollary}

\begin{proof}
    By \eqref{prop:5.6}, $B/\mathfrak{q}$ is integral over $A/\mathfrak{p}$, and both these rings are integral domains. Now use \eqref{prop:5.7}.
\end{proof}

\begin{corollary}\label{cor:5.9}
    Let $A \subseteq B$ be rings, $B$ integral over $A$; let $\mathfrak{q}, \mathfrak{q}'$ be prime ideals of $B$ such that $\mathfrak{q} \subseteq \mathfrak{q}'$ and $\mathfrak{q}^c = \mathfrak{q}'^c = \mathfrak{p}$ say. Then $\mathfrak{q} = \mathfrak{q}'$.
\end{corollary}

\begin{proof}
    By \eqref{prop:5.6}, $B/\mathfrak{q}$ is integral over $A/\mathfrak{p}$. Let $\mathfrak{m}$ be the extension of $\mathfrak{p}$ in $A/\mathfrak{p}$ and let $\mathfrak{n}, \mathfrak{n}'$ be the extensions of $\mathfrak{q}, \mathfrak{q}'$ respectively in $B/\mathfrak{q}$. Then $\mathfrak{m}$ is the maximal ideal of $A/\mathfrak{p}$; $\mathfrak{n} \subseteq \mathfrak{n}'$, and $\mathfrak{n}^c = \mathfrak{n}'^c = \mathfrak{m}$. By \eqref{cor:5.8} it follows that $\mathfrak{n}, \mathfrak{n}'$ are maximal, hence $\mathfrak{n} = \mathfrak{n}'$, hence by \eqref{prop:extension-ideals}\eqref{item 4:prop:extension-ideals}, $\mathfrak{q} = \mathfrak{q}'$.
\end{proof}

\begin{theorem}\label{thm:5.10}\index{lying over}
    Let $A \subseteq B$ be rings, $B$ integral over $A$, and let $\mathfrak{p}$ be a prime ideal of $A$. Then there exists a prime ideal $\mathfrak{q}$ of $B$ such that $\mathfrak{q} \cap A = \mathfrak{p}$.
\end{theorem}
\begin{proof}
    By \eqref{prop:5.6}, $B_{\mathfrak{p}}$ is integral over $A_{\mathfrak{p}}$, and the diagram
    \[\begin{tikzcd}
        A & B \\
        {A_{\mathfrak{p}}} & {B_{\mathfrak{p}}}
        \arrow[from=1-1, to=1-2]
        \arrow["\alpha"', from=1-1, to=2-1]
        \arrow["\beta", from=1-2, to=2-2]
        \arrow[from=2-1, to=2-2]
    \end{tikzcd}\]
    (in which the horizontal arrows are injections) is commutative. Let $\mathfrak{n}$ be a maximal ideal of $B_{\mathfrak{p}}$; then $\mathfrak{m} = \mathfrak{n} \cap A_{\mathfrak{p}}$ is maximal by \eqref{cor:5.8}, hence is the unique maximal ideal of the local ring $A_{\mathfrak{p}}$. If $\mathfrak{q} = \beta^{-1}(\mathfrak{n})$, then $\mathfrak{q}$ is prime and we have $\mathfrak{q} \cap A = \alpha^{-1}(\mathfrak{m}) = \mathfrak{p}$.
\end{proof}

\begin{theorem}[Going-Up Theorem]\label{thm:5.11}\index{Going-Up Theorem}\index{integral extension!going-up}
Let $A \subseteq B$ be rings, $B$ integral over $A$; let $\mathfrak{p}_1 \subseteq \cdots \subseteq \mathfrak{p}_n$ be a chain of prime ideals of $A$ and $\mathfrak{q}_1 \subseteq \cdots \subseteq \mathfrak{q}_m$ ($m < n$) a chain of prime ideals of $B$ such that $\mathfrak{q}_i \cap A = \mathfrak{p}_i$ ($1 \leq i \leq m$). Then the chain $\mathfrak{q}_1 \subseteq \cdots \subseteq \mathfrak{q}_m$ can be extended to a chain $\mathfrak{q}_1 \subseteq \cdots \subseteq \mathfrak{q}_n$ such that $\mathfrak{q}_i \cap A = \mathfrak{p}_i$ for $1 \leq i \leq n$.
\end{theorem}

\begin{proof}
By induction we reduce immediately to the case $m = 1$, $n = 2$. Let $\overline{A} = A/\mathfrak{p}_1$, $\overline{B} = B/\mathfrak{q}_1$; then $\overline{A} \subseteq \overline{B}$, and $\overline{B}$ is integral over $\overline{A}$ by \eqref{prop:5.6}. Hence, by \eqref{thm:5.10}, there exists a prime ideal $\overline{\mathfrak{q}}_2$ of $\overline{B}$ such that $\overline{\mathfrak{q}}_2 \cap \overline{A} = \overline{\mathfrak{p}}_2$, the image of $\mathfrak{p}_2$ in $\overline{A}$. Lift back $\overline{\mathfrak{q}}_2$ to $B$ and we have a prime ideal $\mathfrak{q}_2$ with the required properties.
\end{proof}

\section*{Integrally Closed Integral Domains. The Going-Down Theorem}
\addcontentsline{toc}{section}{eIntegrally Closed Integral Domains. The Going-Down Theorem}

\eqref{prop:5.6} can be sharpened:

\begin{proposition}\label{prop:5.12}
Let $A \subseteq B$ be rings, $C$ the integral closure of $A$ in $B$. Let $S$ be a multiplicatively closed subset of $A$. Then $S^{-1}C$ is the integral closure of $S^{-1}A$ in $S^{-1}B$.
\end{proposition}

\begin{proof}
By \eqref{prop:5.6}, $S^{-1}C$ is integral over $S^{-1}A$. Conversely, if $b/s \in S^{-1}B$ is integral over $S^{-1}A$, then we have an equation of the form
\[
(b/s)^n + (a_1/s_1)(b/s)^{n-1} + \cdots + a_n/s_n = 0
\]
where $a_i \in A$, $s_i \in S$ ($1 \leq i \leq n$). Let $t = s_1 \cdots s_n$ and multiply this equation by $(st)^n$ throughout. Then it becomes an equation of integral dependence for $bt$ over $A$. Hence $bt \in C$ and therefore $b/s = bt/st \in S^{-1}C$.
\end{proof}

An integral domain is said to be \emph{integrally closed} (without qualification) if it is integrally closed in its field of fractions.\index{integrally closed domain} For example, $\Z$ is integrally closed (see Example \eqref{ex:5.0}). The same argument shows that any unique factorization domain is integrally closed. In particular, a polynomial ring $k[x_1, \ldots, x_n]$ over a field is integrally closed.

Integral closure is a local property:

\begin{proposition}\label{prop:5.13}\index{integrally closed!local property}
Let $A$ be an integral domain. Then the following are equivalent:
\begin{enumerate}
    \item[i)] $A$ is integrally closed;
    \item[ii)] $A_{\mathfrak{p}}$ is integrally closed, for each prime ideal $\mathfrak{p}$;
    \item[iii)] $A_{\mathfrak{m}}$ is integrally closed, for each maximal ideal $\mathfrak{m}$.
\end{enumerate}
\end{proposition}
\begin{proof}
    Let $K$ be the field of fractions of $A$, let $C$ be the integral closure of $A$ in $K$, and let $f \colon A \to C$ be the identity mapping of $A$ into $C$. Then $A$ is integrally closed $\iff$ $f$ is surjective, and by \eqref{prop:5.12}, $A_{\mathfrak{p}}$ (resp.~$A_{\mathfrak{m}}$) is integrally closed $\iff$ $f_{\mathfrak{p}}$ (resp.~$f_{\mathfrak{m}}$) is surjective. Now use \eqref{prop:injective-surjective-local}.
\end{proof}

Let $A \subseteq B$ be rings and let $\mathfrak{a}$ be an ideal of $A$. An element of $B$ is said to be \emph{integral over $\mathfrak{a}$} if it satisfies an equation of integral dependence over $A$ in which all the coefficients lie in $\mathfrak{a}$. The \emph{integral closure} of $\mathfrak{a}$ in $B$ is the set of all elements of $B$ which are integral over $\mathfrak{a}$.

\begin{lemma}\label{lem:5.14}
Let $C$ be the integral closure of $A$ in $B$ and let $\mathfrak{a}^e$ denote the extension of $\mathfrak{a}$ in $C$. Then the integral closure of $\mathfrak{a}$ in $B$ is the radical of $\mathfrak{a}^e$ (and is therefore closed under addition and multiplication).
\end{lemma}

\begin{proof}
If $x \in B$ is integral over $\mathfrak{a}$, we have an equation of the form
\[
x^n + a_1 x^{n-1} + \cdots + a_n = 0
\]
with $a_1, \ldots, a_n$ in $\mathfrak{a}$. Hence $x \in C$ and $x^n \in \mathfrak{a}^e$, that is $x \in r(\mathfrak{a}^e)$. Conversely, if $x \in r(\mathfrak{a}^e)$ then $x^n = \sum a_i x_i$ for some $n > 0$, where the $a_i$ are elements of $\mathfrak{a}$ and the $x_i$ are elements of $C$. Since each $x_i$ is integral over $A$ it follows from \eqref{cor:5.2} that $M = A[x_1, \ldots, x_n]$ is a finitely generated $A$-module, and we have $x^n M \subseteq \mathfrak{a} M$. Hence by \eqref{prop:module-endomorphism-equation} (taking $\phi$ there to be multiplication by $x^n$) we see that $x^n$ is integral over $\mathfrak{a}$, hence $x$ is integral over $\mathfrak{a}$.
\end{proof}

\begin{proposition}\label{prop:5.15}
    Let $A \subseteq B$ be integral domains, $A$ integrally closed, and let $x \in B$ be integral over an ideal $\mathfrak{a}$ of $A$. Then $x$ is algebraic over the field of fractions $K$ of $A$, and if its minimal polynomial over $K$ is $t^n + a_1 t^{n-1} + \cdots + a_n$, then $a_1, \ldots, a_n$ lie in $r(\mathfrak{a})$.
\end{proposition}

\begin{proof}
Clearly $x$ is algebraic over $K$. Let $L$ be an extension field of $K$ which contains all the conjugates $x_1, \ldots, x_n$ of $x$. Each $x_i$ satisfies the same equation of integral dependence as $x$ does, hence each $x_i$ is integral over $\mathfrak{a}$. The coefficients of the minimal polynomial of $x$ over $K$ are polynomials in the $x_i$, hence by \eqref{lem:5.14} are integral over $\mathfrak{a}$. Since $A$ is integrally closed, they must lie in $r(\mathfrak{a})$, by \eqref{lem:5.14} again.
\end{proof}

\begin{theorem}[Going-Down Theorem]\label{thm:5.16}\index{Going-Down Theorem}\index{integral extension!going-down}
Let $A \subseteq B$ be integral domains, $A$ integrally closed, $B$ integral over $A$. Let $\mathfrak{p}_1 \supseteq \cdots \supseteq \mathfrak{p}_n$ be a chain of prime ideals of $A$, and let $\mathfrak{q}_1 \supseteq \cdots \supseteq \mathfrak{q}_m$ ($m < n$) be a chain of prime ideals of $B$ such that $\mathfrak{q}_i \cap A = \mathfrak{p}_i$ ($1 \leq i \leq m$). Then the chain $\mathfrak{q}_1 \supseteq \cdots \supseteq \mathfrak{q}_m$ can be extended to a chain $\mathfrak{q}_1 \supseteq \cdots \supseteq \mathfrak{q}_n$ such that $\mathfrak{q}_i \cap A = \mathfrak{p}_i$ ($1 \leq i \leq n$).
\end{theorem}

\begin{proof}
    As in \eqref{thm:5.11} we reduce immediately to the case $m = 1$, $n = 2$. Then we have to show that $\mathfrak{p}_2$ is the contraction of a prime ideal in the ring $B_{\mathfrak{q}_1}$, or equivalently \eqref{prop:contracted-ideal} that $B_{\mathfrak{q}_1} \mathfrak{p}_2 \cap A = \mathfrak{p}_2$.

    Every $x \in B_{\mathfrak{q}_1} \mathfrak{p}_2$ is of the form $y/s$, where $y \in B \mathfrak{p}_2$ and $s \in B - \mathfrak{q}_1$. By \eqref{lem:5.14}, $y$ is integral over $\mathfrak{p}_2$, and hence by \eqref{prop:5.15} its minimal equation over $K$, the field of fractions of $A$, is of the form
    \begin{equation}\label{eq:5.16-1}
    y^r + u_1 y^{r-1} + \cdots + u_r = 0 \tag{1}
    \end{equation}
    with $u_1, \ldots, u_r$ in $\mathfrak{p}_2$.

    Now suppose that $x \in B_{\mathfrak{q}_1} \mathfrak{p}_2 \cap A$. Then $s = y x^{-1}$ with $x^{-1} \in K$, so that the minimal equation for $s$ over $K$ is obtained by dividing \eqref{eq:5.16-1} by $x^r$, and is therefore, say,
    \begin{equation}\label{eq:5.16-2}
    s^r + v_1 s^{r-1} + \cdots + v_r = 0 \tag{2}
    \end{equation}
    where $v_i = u_i / x^i$. Consequently
    \begin{equation}\label{eq:5.16-3}
    x^i v_i = u_i \in \mathfrak{p}_2 \qquad (1 \leq i \leq r) \tag{3}.
    \end{equation}

    But $s$ is integral over $A$, hence by \eqref{prop:5.15} (with $\mathfrak{a} = (1)$) each $v_i$ is in $A$. Suppose $x \notin \mathfrak{p}_2$. Then \eqref{eq:5.16-3} shows that each $v_i \in \mathfrak{p}_2$, hence \eqref{eq:5.16-2} shows that $s^r \in B \mathfrak{p}_2 \subseteq B \mathfrak{p}_1 \subseteq \mathfrak{q}_1$, and therefore $s \in \mathfrak{q}_1$, which is a contradiction. Hence $x \in \mathfrak{p}_2$ and therefore $B_{\mathfrak{q}_1} \mathfrak{p}_2 \cap A = \mathfrak{p}_2$ as required.
\end{proof}

The proof of the next proposition assumes some standard facts from field theory.

\begin{proposition}\label{prop:5.17}
    Let $A$ be an integrally closed domain, $K$ its field of fractions, $L$ a finite separable algebraic extension of $K$, $B$ the integral closure of $A$ in $L$. Then there exists a basis $v_1, \ldots, v_n$ of $L$ over $K$ such that $B \subseteq \sum_{i=1}^n A v_i$.
\end{proposition}

\begin{proof}
If $v$ is any element of $L$, then $v$ is algebraic over $K$ and therefore satisfies an equation of the form
\[
a_0 v^r + a_1 v^{r-1} + \cdots + a_r = 0 \qquad (a_i \in A).
\]
Multiplying this equation by $a_0^{r-1}$, we see that $a_0 v = u$ is integral over $A$, and hence is in $B$. Thus, given any basis of $L$ over $K$ we may multiply the basis elements by suitable elements of $A$ to get a basis $u_1, \ldots, u_n$ such that each $u_i \in B$.

Let $T$ denote trace (from $L$ to $K$). Since $L/K$ is separable, the bilinear form $(x, y) \mapsto T(xy)$ on $L$ (considered as a vector space over $K$) is non-degenerate, and hence we have a dual basis $v_1, \ldots, v_n$ of $L$ over $K$, defined by $T(u_i v_j) = \delta_{ij}$. Let $x \in B$, say $x = \sum x_i v_i$ ($x_i \in K$). We have $x u_i \in B$ (since $u_i \in B$) and therefore $T(x u_i) \in A$ by \eqref{prop:5.15} (for the trace of an element is a multiple of one of the coefficients in the minimal polynomial). But $T(x u_i) = \sum_j T(x_j u_i v_j) = \sum_j x_j T(u_i v_j) = \sum_j x_j \delta_{ij} = x_i$, hence $x_i \in A$. Consequently $B \subseteq \sum_i A v_i$.
\end{proof}

\section*{Valuation Rings}
\addcontentsline{toc}{section}{Valuation Rings}

Let $B$ be an integral domain, $K$ its field of fractions. $B$ is a \emph{valuation ring}\index{valuation ring} of $K$ if, for each $x \neq 0$, either $x \in B$ or $x^{-1} \in B$ (or both).\index{valuation ring}

\begin{proposition}\label{prop:5.18}\index{valuation ring!properties}
\begin{enumerate}
    \item[i)] $B$ is a local ring.
    \item[ii)] If $B'$ is a ring such that $B \subseteq B' \subseteq K$, then $B'$ is a valuation ring of $K$.
    \item[iii)] $B$ is integrally closed (in $K$).
\end{enumerate}
\end{proposition}

\begin{proof}
    i) Let $\mathfrak{m}$ be the set of non-units of $B$, so that $x \in \mathfrak{m} \iff$ either $x = 0$ or $x^{-1} \notin B$. If $a \in B$ and $x \in \mathfrak{m}$ we have $ax \in \mathfrak{m}$, for otherwise $(ax)^{-1} \in B$ and therefore $x^{-1} = a \cdot (ax)^{-1} \in B$. Next let $x, y$ be non-zero elements of $\mathfrak{m}$. Then either $x y^{-1} \in B$ or $x^{-1} y \in B$. If $x y^{-1} \in B$ then $x + y = (1 + x y^{-1})y \in B \mathfrak{m} \subseteq \mathfrak{m}$, and similarly if $x^{-1} y \in B$. Hence $\mathfrak{m}$ is an ideal and therefore $B$ is a local ring by \eqref{prop:local-ring-characterization}.

    ii) Clear from the definitions.

    iii) Let $x \in K$ be integral over $B$. Then we have, say,
    \[
    x^n + b_1 x^{n-1} + \cdots + b_n = 0
    \]
    with the $b_i \in B$. If $x \in B$ there is nothing to prove. If not, then $x^{-1} \in B$, hence $x = -(b_1 + b_2 x^{-1} + \cdots + b_n x^{1-n}) \in B$.
\end{proof}

Let $K$ be a field, $\Omega$ an algebraically closed field. Let $\Sigma$ be the set of all pairs $(A, f)$, where $A$ is a subring of $K$ and $f$ is a homomorphism of $A$ into $\Omega$. We partially order the set $\Sigma$ as follows:
\[
(A, f) \leq (A', f') \quad \iff \quad A \subseteq A' \text{ and } f'|_A = f.
\]
The conditions of Zorn's lemma are clearly satisfied and therefore the set $\Sigma$ has at least one maximal element.

Let $(B, g)$ be a maximal element of $\Sigma$. We want to prove that $B$ is a valuation ring of $K$. The first step in the proof is

\begin{lemma}\label{lem:5.19}
$B$ is a local ring and $\mathfrak{m} = \Ker(g)$ is its maximal ideal.
\end{lemma}

\begin{proof}
Since $g(B)$ is a subring of a field and therefore an integral domain, the ideal $\mathfrak{m} = \Ker(g)$ is prime. We can extend $g$ to a homomorphism $\bar{g} \colon B_{\mathfrak{m}} \to \Omega$ by putting $\bar{g}(b/s) = g(b)/g(s)$ for all $b \in B$ and all $s \in B - \mathfrak{m}$, since $g(s)$ will not be zero. Since the pair $(B, g)$ is maximal it follows that $B = B_{\mathfrak{m}}$, hence $B$ is a local ring and $\mathfrak{m}$ is its maximal ideal.
\end{proof}

\begin{lemma}\label{lem:5.20}
Let $x$ be a non-zero element of $K$. Let $B[x]$ be the subring of $K$ generated by $x$ over $B$, and let $\mathfrak{m}[x]$ be the extension of $\mathfrak{m}$ in $B[x]$. Then either $\mathfrak{m}[x] \neq B[x]$ or $\mathfrak{m}[x^{-1}] \neq B[x^{-1}]$.
\end{lemma}

\begin{proof}
Suppose that $\mathfrak{m}[x] = B[x]$ and $\mathfrak{m}[x^{-1}] = B[x^{-1}]$. Then we shall have equations
\begin{align}
u_0 + u_1 x + \cdots + u_m x^m &= 1 \qquad (u_i \in \mathfrak{m})  \tag{1}\label{eq:5.20-1} \\
v_0 + v_1 x^{-1} + \cdots + v_n x^{-n} &= 1 \qquad (v_j \in \mathfrak{m}) \tag{2} \label{eq:5.20-2}
\end{align}
in which we may assume that the degrees $m, n$ are as small as possible. Suppose that $m \geq n$, and multiply \eqref{eq:5.20-2} through by $x^n$:
\begin{equation}\label{eq:5.20-3}
v_0 x^n + v_1 x^{n-1} + \cdots + v_n = x^n. \tag{3}
\end{equation}
Since $v_0 \in \mathfrak{m}$, it follows from \eqref{lem:5.19} that $1 - v_0$ is a unit in $B$, and \eqref{eq:5.20-3} may therefore be written in the form
\[
x^n = w_1 x^{n-1} + \cdots + w_n \qquad (w_j \in \mathfrak{m}).
\]
Hence we can replace $x^m$ in \eqref{eq:5.20-1} by $w_1 x^{m-1} + \cdots + w_n x^{m-n}$, and this contradicts the minimality of the exponent $m$.
\end{proof}

\begin{theorem}\label{thm:5.21}\index{valuation ring!existence}
Let $(B, g)$ be a maximal element of $\Sigma$. Then $B$ is a valuation ring of the field $K$.
\end{theorem}

\begin{proof}
We have to show that if $x \neq 0$ is an element of $K$, then either $x \in B$ or $x^{-1} \in B$. By \eqref{lem:5.20} we may as well assume that $\mathfrak{m}[x]$ is not the unit ideal of the ring $B' = B[x]$. Then $\mathfrak{m}[x]$ is contained in a maximal ideal $\mathfrak{m}'$ of $B'$, and we have $\mathfrak{m}' \cap B = \mathfrak{m}$ (because $\mathfrak{m}' \cap B$ is a proper ideal of $B$ and contains $\mathfrak{m}$). Hence the embedding of $B$ in $B'$ induces an embedding of the field $k = B/\mathfrak{m}$ in the field $k' = B'/\mathfrak{m}'$; also $k' = k[\bar{x}]$ where $\bar{x}$ is the image of $x$ in $k'$, hence $\bar{x}$ is algebraic over $k$, and therefore $k'$ is a finite algebraic extension of $k$.

Now the homomorphism $g$ induces an embedding $\bar{g}$ of $k$ into $\Omega$, since by \eqref{lem:5.19}, $\mathfrak{m}$ is the kernel of $g$. Since $\Omega$ is algebraically closed, $\bar{g}$ can be extended to an embedding $\bar{g}'$ of $k'$ into $\Omega$. Composing $\bar{g}'$ with the natural homomorphism $B' \to k'$, we have, say, $g' \colon B' \to \Omega$ which extends $g$. Since the pair $(B, g)$ is maximal, it follows that $B' = B$ and therefore $x \in B$.
\end{proof}

\begin{corollary}\label{cor:5.22}\index{integral closure!as intersection of valuation rings}
Let $A$ be a subring of a field $K$. Then the integral closure $\overline{A}$ of $A$ in $K$ is the intersection of all the valuation rings of $K$ which contain $A$.
\end{corollary}

\begin{proof}
Let $B$ be a valuation ring of $K$ such that $A \subseteq B$. Since $B$ is integrally closed, by \eqref{prop:5.18}(iii), it follows that $\overline{A} \subseteq B$.

Conversely, let $x \notin \overline{A}$. Then $x$ is not in the ring $A' = A[x^{-1}]$. Hence $x^{-1}$ is a non-unit in $A'$ and is therefore contained in a maximal ideal $\mathfrak{m}'$ of $A'$. Let $\Omega$ be an algebraic closure of the field $k' = A'/\mathfrak{m}'$. Then the restriction to $A$ of the natural homomorphism $A' \to k'$ defines a homomorphism of $A$ into $\Omega$. By \eqref{thm:5.21} this can be extended to some valuation ring $B \supseteq A$. Since $x^{-1}$ maps to zero, it follows that $x \notin B$.
\end{proof}

\begin{proposition}\label{prop:5.23}
Let $A \subseteq B$ be integral domains, $B$ finitely generated over $A$. Let $v$ be a non-zero element of $B$. Then there exists $u \neq 0$ in $A$ with the following property: any homomorphism $f$ of $A$ into an algebraically closed field $\Omega$ such that $f(u) \neq 0$ can be extended to a homomorphism $g$ of $B$ into $\Omega$ such that $g(v) \neq 0$.
\end{proposition}

\begin{proof}
By induction on the number of generators of $B$ over $A$ we reduce immediately to the case where $B$ is generated over $A$ by a single element $x$.

i) Suppose $x$ is transcendental over $A$, i.e., that no non-zero polynomial with coefficients in $A$ has $x$ as a root. Let $v = a_0 x^n + a_1 x^{n-1} + \cdots + a_n$, and take $u = a_0$. Then if $f \colon A \to \Omega$ is such that $f(u) \neq 0$, there exists $\xi \in \Omega$ such that $f(a_0) \xi^n + f(a_1) \xi^{n-1} + \cdots + f(a_n) \neq 0$, because $\Omega$ is infinite. Define $g \colon B \to \Omega$ extending $f$ by putting $g(x) = \xi$. Then $g(v) \neq 0$, as required.

ii) Now suppose $x$ is algebraic over $A$ (i.e.~over the field of fractions of $A$). Then so is $v^{-1}$, because $v$ is a polynomial in $x$. Hence we have equations of the form
\begin{align}
a_0 x^m + a_1 x^{m-1} + \cdots + a_m &= 0 \qquad (a_i \in A) \label{eq:5.23-1}\tag{1} \\
a'_0 v^{-n} + a'_1 v^{1-n} + \cdots + a'_n &= 0 \qquad (a'_i \in A). \label{eq:5.23-2}\tag{2}
\end{align}
Let $u = a_0 a'_0$, and let $f \colon A \to \Omega$ be such that $f(u) \neq 0$. Then $f$ can be extended, first to a homomorphism $h \colon A[u^{-1}] \to \Omega$ (with $h(u^{-1}) = f(u)^{-1}$), and then by \eqref{thm:5.21} to a homomorphism $h \colon C \to \Omega$, where $C$ is a valuation ring containing $A[u^{-1}]$. From \eqref{eq:5.23-1}, $x$ is integral over $A[u^{-1}]$, hence by \eqref{cor:5.22}, $x \in C$, so that $C$ contains $B$, and in particular $v \in C$. On the other hand, from \eqref{eq:5.23-2}, $v^{-1}$ is integral over $A[u^{-1}]$, and therefore by \eqref{cor:5.22} again is in $C$. Therefore $v$ is a unit in $C$, and hence $h(v) \neq 0$. Now take $g$ to be the restriction of $h$ to $B$.
\end{proof}

\begin{corollary}\label{cor:5.24}\index{Nullstellensatz!weak form}
Let $k$ be a field and $B$ a finitely generated $k$-algebra. If $B$ is a field then it is a finite algebraic extension of $k$.
\end{corollary}

\begin{proof}
Take $A = k$, $v = 1$ and $\Omega =$ algebraic closure of $k$.
\end{proof}

\eqref{cor:5.24} is one form of Hilbert's Nullstellensatz. For another proof, see \eqref{prop:nullstellensatz-weak}.

%=============================================================================
% Chapter 5 Exercises - With hints from original book (35 exercises)
%=============================================================================
\section*{Exercises}
\addcontentsline{toc}{section}{Exercises}

\begin{enumerate}
    \item \label{exer1-chapter5}Let $f \colon A \to B$ be an integral homomorphism of rings. Show that $f^* \colon \Spec(B) \to \Spec(A)$ is a closed mapping, i.e., that it maps closed sets to closed sets. (This is a geometrical equivalent of \eqref{thm:5.10}.)

    \item \label{exer2-chapter5}Let $A$ be a subring of a ring $B$ such that $B$ is integral over $A$, and let $f \colon A \to \Omega$ be a homomorphism of $A$ into an algebraically closed field $\Omega$. Show that $f$ can be extended to a homomorphism of $B$ into $\Omega$.\\
    \textit{Hint:} Use \eqref{thm:5.10}.

    \item \label{exer3-chapter5}Let $f \colon B \to B'$ be a homomorphism of $A$-algebras, and let $C$ be an $A$-algebra. If $f$ is integral, prove that $f \otimes 1 \colon B \otimes_A C \to B' \otimes_A C$ is integral. (This includes (5.6)(ii) as a special case.)

    \item \label{exer4-chapter5}Let $A$ be a subring of a ring $B$ such that $B$ is integral over $A$. Let $\mathfrak{n}$ be a maximal ideal of $B$ and let $\mathfrak{m} = \mathfrak{n} \cap A$ be the corresponding maximal ideal of $A$. Is $B_{\mathfrak{n}}$ necessarily integral over $A_{\mathfrak{m}}$?\\
    \textit{Hint:} Consider the subring $k[x^2 - 1]$ of $k[x]$, where $k$ is a field, and let $\mathfrak{n} = (x - 1)$. Can the element $1/(x + 1)$ be integral over $k[x^2 - 1]_{(x^2-1)}$?

    \item \label{exer5-chapter5}Let $A \subseteq B$ be rings, $B$ integral over $A$.
    \begin{enumerate}[i)]
        \item If $x \in A$ is a unit in $B$ then it is a unit in $A$.
        \item The Jacobson radical of $A$ is the contraction of the Jacobson radical of $B$.
    \end{enumerate}

    \item \label{exer6-chapter5}Let $B_1, \ldots, B_n$ be integral $A$-algebras. Show that $\prod_{i=1}^n B_i$ is an integral $A$-algebra.

    \item \label{exer7-chapter5}Let $A$ be a subring of a ring $B$, such that the set $B \setminus A$ is closed under multiplication. Show that $A$ is integrally closed in $B$.

    \item \label{exer8-chapter5} \begin{enumerate}[i)]
        \item Let $A$ be a subring of an integral domain $B$, and let $C$ be the integral closure of $A$ in $B$. Let $f, g$ be monic polynomials in $B[x]$ such that $fg \in C[x]$. Then $f, g \in C[x]$.\\
        \textit{Hint:} Take a field containing $B$ in which the polynomials $f, g$ split into linear factors: say $f = \prod (x - \xi_i)$, $g = \prod (x - \eta_j)$. Each $\xi_i$ and each $\eta_j$ is a root of $fg$, hence is integral over $C$. Hence the coefficients of $f$ and $g$ are integral over $C$.
        \item Prove the same result without assuming that $B$ (or $A$) is an integral domain.
    \end{enumerate}

    \item \label{exer9-chapter5}Let $A$ be a subring of a ring $B$ and let $C$ be the integral closure of $A$ in $B$. Prove that $C[x]$ is the integral closure of $A[x]$ in $B[x]$.\\
    \textit{Hint:} If $f \in B[x]$ is integral over $A[x]$, then
    \[
    f^m + g_1 f^{m-1} + \cdots + g_m = 0 \quad (g_j \in A[x]).
    \]
    Let $r$ be an integer larger than $m$ and the degrees of $g_1, \ldots, g_m$, and let $f_1 = f - x^r$ so that
    \[
    f_1^m + h_1 f_1^{m-1} + \cdots + h_m = 0,
    \]
    where $h_m = (x^r)^m + g_1(x^r)^{m-1} + \cdots + g_m \in A[x]$. Now apply Exercise \ref{exer8-chapter5} to the polynomials $-f_1$ and $f_1^{m-1} + h_1 f_1^{m-2} + \cdots + h_{m-1}$.

    \item \label{exer10-chapter5} A ring homomorphism $f \colon A \to B$ is said to have the \emph{going-up property} (resp. the \emph{going-down property}) if the conclusion of the going-up theorem \eqref{thm:5.11} (resp. the going-down theorem \eqref{thm:5.16}) holds for $B$ and its subring $f(A)$.\\
    Let $f^* \colon \Spec(B) \to \Spec(A)$ be the mapping associated with $f$.
    \begin{enumerate}[i)]
        \item Consider the following three statements:
        \begin{itemize}
            \item[(a)] $f^*$ is a closed mapping.
            \item[(b)] $f$ has the going-up property.
            \item[(c)] Let $\mathfrak{q}$ be any prime ideal of $B$ and let $\mathfrak{p} = \mathfrak{q}^c$. Then $f^* \colon \Spec(B/\mathfrak{q}) \to \Spec(A/\mathfrak{p})$ is surjective.
        \end{itemize}
        Prove that (a) $\implies$ (b) $\iff$ (c). (See also Chapter \ref{chapter6}, Exercise \ref{exer11-chapter6}.)
        \item Consider the following three statements:
        \begin{itemize}
            \item[(a')] $f^*$ is an open mapping.
            \item[(b')] $f$ has the going-down property.
            \item[(c')] For any prime ideal $\mathfrak{q}$ of $B$, if $\mathfrak{p} = \mathfrak{q}^c$, then $f^* \colon \Spec(B_{\mathfrak{q}}) \to \Spec(A_{\mathfrak{p}})$ is surjective.
        \end{itemize}
        Prove that (a') $\implies$ (b') $\iff$ (c'). (See also Chapter \ref{chapter7}, Exercise \ref{exer23-chapter7}.)
    \end{enumerate}
    \textit{Hint for (ii):} To prove that (a') $\implies$ (c'), observe that $B_{\mathfrak{q}}$ is the direct limit of the rings $B_t$ where $t \in B \setminus \mathfrak{q}$; hence, by Chapter \ref{chapter3}, Exercise \ref{exer26-chapter3}, we have $f^*(\Spec(B_{\mathfrak{q}})) = \bigcap f_t^*(\Spec(B_t)) = \bigcap f_t^*(Y_t)$. Since $Y_t$ is an open neighborhood of $\mathfrak{q}$ in $Y$, and since $f^*$ is open, it follows that $f^*(Y_t)$ is an open neighborhood of $\mathfrak{p}$ in $X$ and therefore contains $\Spec(A_{\mathfrak{p}})$.

    \item \label{exer11-chapter5} Let $f \colon A \to B$ be a flat homomorphism of rings. Then $f$ has the going-down property.
    
    \textit{Hint:} Chapter \ref{chapter3}, Exercise \ref{exer18-chapter3}.

    \item \label{exer12-chapter5} Let $G$ be a finite group of automorphisms of a ring $A$, and let $A^G$ denote the subring of $G$-invariants, that is of all $x \in A$ such that $\sigma(x) = x$ for all $\sigma \in G$. Prove that $A$ is integral over $A^G$.

    \textit{Hint:} If $x \in A$, observe that $x$ is a root of the polynomial $\prod_{\sigma \in G} (t - \sigma(x))$.\\
    Let $S$ be a multiplicatively closed subset of $A$ such that $\sigma(S) \subseteq S$ for all $\sigma \in G$, and let $S^G = S \cap A^G$. Show that the action of $G$ on $A$ extends to an action on $S^{-1}A$, and that $(S^G)^{-1}A^G \cong (S^{-1}A)^G$.

    \item \label{exer13-chapter5}In the situation of Exercise \ref{exer12-chapter5}, let $\mathfrak{p}$ be a prime ideal of $A^G$, and let $P$ be the set of prime ideals of $A$ whose contraction is $\mathfrak{p}$. Show that $G$ acts transitively on $P$. In particular, $P$ is finite.

    \textit{Hint:} Let $\mathfrak{P}_1, \mathfrak{P}_2 \in P$ and let $x \in \mathfrak{P}_1$. Then $\prod_{\sigma} \sigma(x) \in \mathfrak{P}_1 \cap A^G = \mathfrak{p} \subseteq \mathfrak{P}_2$, hence $\sigma(x) \in \mathfrak{P}_2$ for some $\sigma \in G$. Deduce that $\mathfrak{P}_1$ is contained in $\bigcup_{\sigma \in G} \sigma(\mathfrak{P}_2)$, and then apply \eqref{pro:prime avoidance} and \eqref{cor:5.9}.

    \item \label{exer14-chapter5} Let $A$ be an integrally closed domain, $K$ its field of fractions and $L$ a finite normal separable extension of $K$. Let $G$ be the Galois group of $L$ over $K$ and let $B$ be the integral closure of $A$ in $L$. Show that $\sigma(B) = B$ for all $\sigma \in G$, and that $A = B^G$.

    \item \label{exer15-chapter5} Let $A, K$ be as in Exercise \ref{exer14-chapter5}, let $L$ be any finite extension field of $K$, and let $B$ be the integral closure of $A$ in $L$. Show that, if $\mathfrak{p}$ is any prime ideal of $A$, then the set of prime ideals $\mathfrak{q}$ of $B$ which contract to $\mathfrak{p}$ is finite (in other words, that $\Spec(B) \to \Spec(A)$ has finite fibers).

    \textit{Hint:} Reduce to the two cases (a) $L$ separable over $K$ and (b) $L$ purely inseparable over $K$. In case (a), embed $L$ in a finite normal separable extension of $K$, and use Exercises \ref{exer13-chapter5} and \ref{exer14-chapter5}. In case (b), if $\mathfrak{q}$ is a prime ideal of $B$ such that $\mathfrak{q} \cap A = \mathfrak{p}$, show that $\mathfrak{q}$ is the set of all $x \in B$ such that $x^{p^m} \in \mathfrak{p}$ for some $m \geq 0$, where $p$ is the characteristic of $K$, and hence that $\Spec(B) \to \Spec(A)$ is bijective in this case.

    \emph{Noether's normalization lemma.}

    \item \label{exer16-chapter5} Let $k$ be a field and let $A \neq 0$ be a finitely generated $k$-algebra. Then there exist elements $y_1, \ldots, y_r \in A$ which are algebraically independent over $k$ and such that $A$ is integral over $k[y_1, \ldots, y_r]$.

    \quad\, We shall assume that $k$ is infinite. (The result is still true if $k$ is finite, but a different proof is needed.) Let $x_1, \ldots, x_n$ generate $A$ as a $k$-algebra. We can renumber the $x_i$ so that $x_1, \ldots, x_r$ are algebraically independent over $k$ and each of $x_{r+1}, \ldots, x_n$ is algebraic over $k[x_1, \ldots, x_r]$. Now proceed by induction on $n$. If $n = r$ there is nothing to do, so suppose $n > r$ and the result true for $n-1$ generators. The generator $x_n$ is algebraic over $k[x_1, \ldots, x_{n-1}]$, hence there exists a polynomial $f \neq 0$ in $n$ variables such that $f(x_1, \ldots, x_{n-1}, x_n) = 0$. Let $F$ be the homogeneous part of highest degree in $f$. Since $k$ is infinite, there exist $\lambda_1, \ldots, \lambda_{n-1} \in k$ such that $F(\lambda_1, \ldots, \lambda_{n-1}, 1) \neq 0$. Put $x_i' = x_i - \lambda_i x_n$ ($1 \leq i \leq n-1$). Show that $x_n$ is integral over the ring $A' = k[x_1', \ldots, x_{n-1}']$, and hence that $A$ is integral over $A'$. Then apply the inductive hypothesis to $A'$ to complete the proof.

    \quad\, From the proof it follows that $y_1, \ldots, y_r$ may be chosen to be linear combinations of $x_1, \ldots, x_n$. This has the following geometrical interpretation: if $k$ is algebraically closed and $X$ is an affine algebraic variety in $k^n$ with coordinate ring $A \neq 0$, then there exists a linear subspace $L$ of dimension $r$ in $k^n$ and a linear mapping of $k^n$ onto $L$ which maps $X$ onto $L$. \textit{Hint:} Use Exercise \ref{exer2-chapter5}.

    \emph{Nullstellensatz (weak form).}

    \item \label{exer17-chapter5} Let $X$ be an affine algebraic variety in $k^n$, where $k$ is an algebraically closed field, and let $I(X)$ be the ideal of $X$ in the polynomial ring $k[t_1, \ldots, t_n]$ (Chapter \ref{chapter1}, Exercise \ref{exer27-chapter1}). If $I(X) \neq (1)$ then $X$ is not empty.

    \textit{Hint:} Let $A = k[t_1, \ldots, t_n]/I(X)$ be the coordinate ring of $X$. Then $A \neq 0$, hence by Exercise \ref{exer16-chapter5} there exists a linear subspace $L$ of dimension $\geq 0$ in $k^n$ and a mapping of $X$ onto $L$. Hence $X \neq \varnothing$.

    \quad\, Deduce that every maximal ideal in the ring $k[t_1, \ldots, t_n]$ is of the form $(t_1 - a_1, \ldots, t_n - a_n)$ where $a_i \in k$.

    \item \label{exer18-chapter5}Let $k$ be a field and let $B$ be a finitely generated $k$-algebra. Suppose that $B$ is a field. Then $B$ is a finite algebraic extension of $k$. (This is another version of Hilbert's Nullstellensatz. The following proof is due to Zariski. For other proofs, see \eqref{cor:5.24}, \eqref{prop:nullstellensatz-weak}.)
    
    \quad\, Let $x_1, \ldots, x_n$ generate $B$ as a $k$-algebra. The proof is by induction on $n$. If $n = 1$ the result is clearly true, so assume $n > 1$. Let $A = k[x_1]$ and let $K = k(x_1)$ be the field of fractions of $A$. By the inductive hypothesis, $B$ is a finite algebraic extension of $K$, hence each of $x_2, \ldots, x_n$ satisfies a monic polynomial equation with coefficients in $K$, i.e., coefficients of the form $a/b$ where $a$ and $b$ are in $A$. If $f$ is the product of the denominators of all these coefficients, then each of $x_2, \ldots, x_n$ is integral over $A_f$. Hence $B$ and therefore $K$ is integral over $A_f$.

    \quad\, Suppose $x_1$ is transcendental over $k$. Then $A$ is integrally closed, because it is a unique factorization domain. Hence $A_f$ is integrally closed \eqref{prop:5.12}, and therefore $A_f = K$, which is clearly absurd. Hence $x_1$ is algebraic over $k$, hence $K$ (and therefore $B$) is a finite extension of $k$.

    \item \label{exer19-chapter5}Deduce the result of Exercise \ref{exer17-chapter5} from Exercise \ref{exer18-chapter5}.

    \item \label{exer20-chapter5}Let $A$ be a subring of an integral domain $B$ such that $B$ is finitely generated over $A$. Show that there exists $s \neq 0$ in $A$ and elements $y_1, \ldots, y_n$ in $B$, algebraically independent over $A$ and such that $B_s$ is integral over $B'_s$, where $B' = A[y_1, \ldots, y_n]$.
    
    \textit{Hint:} Let $S = A \setminus \{0\}$ and let $K = S^{-1}A$, the field of fractions of $A$. Then $S^{-1}B$ is a finitely generated $K$-algebra and therefore by the normalization lemma (Exercise \ref{exer16-chapter5}) there exist $x_1, \ldots, x_n$ in $S^{-1}B$, algebraically independent over $K$ and such that $S^{-1}B$ is integral over $K[x_1, \ldots, x_n]$. Let $z_1, \ldots, z_m$ generate $B$ as an $A$-algebra. Then each $z_j$ (regarded as an element of $S^{-1}B$) is integral over $K[x_1, \ldots, x_n]$. By writing an equation of integral dependence for each $z_j$, show that there exists $s \in S$ such that $x_i = y_i/s$ ($1 \leq i \leq n$) with $y_i \in B$, and such that each $sz_j$ is integral over $B'$. Deduce that this $s$ satisfies the conditions stated.

    \item \label{exer21-chapter5}Let $A, B$ be as in Exercise \ref{exer20-chapter5}. Show that there exists $s \neq 0$ in $A$ such that, if $\Omega$ is an algebraically closed field and $f \colon A \to \Omega$ is a homomorphism for which $f(s) \neq 0$, then $f$ can be extended to a homomorphism $B \to \Omega$.\\
    \textit{Hint:} With the notation of Exercise \ref{exer20-chapter5}, $f$ can be extended first of all to $B'$, for example by mapping each $y_i$ to $0$; then to $B'_s$ (because $f(s) \neq 0$), and finally to $B_s$ (by Exercise \ref{exer2-chapter5} because $B_s$ is integral over $B'_s$).

    \item \label{exer22-chapter5}Let $A, B$ be as in Exercise \ref{exer20-chapter5}. If the Jacobson radical of $A$ is zero, then so is the Jacobson radical of $B$.
    
    \textit{Hint:} Let $v \neq 0$ be an element of $B$. We have to show that there is a maximal ideal of $B$ which does not contain $v$. By applying Exercise \ref{exer21-chapter5} to the ring $B_v$ and its subring $A$, we obtain an element $s \neq 0$ in $A$. Let $\mathfrak{m}$ be a maximal ideal of $A$ such that $s \notin \mathfrak{m}$, and let $k = A/\mathfrak{m}$. Then the canonical mapping $A \to k$ extends to a homomorphism $g$ of $B_v$ into an algebraic closure $\Omega$ of $k$. Show that $g(v) \neq 0$ and that $\Ker(g) \cap B$ is a maximal ideal of $B$.

    \item \label{exer23-chapter5}Let $A$ be a ring. Show that the following are equivalent:
    \begin{enumerate}[i)]
        \item Every prime ideal in $A$ is an intersection of maximal ideals.
        \item In every homomorphic image of $A$ the nilradical is equal to the Jacobson radical.
        \item Every prime ideal in $A$ which is not maximal is equal to the intersection of the prime ideals which contain it strictly.
    \end{enumerate}
    \textit{Hint:} The only hard part is iii) $\implies$ ii). Suppose ii) false, then there is a prime ideal which is not an intersection of maximal ideals. Passing to the quotient ring, we may assume that $A$ is an integral domain whose Jacobson radical $\mathfrak{R}$ is not zero. Let $f$ be a non-zero element of $\mathfrak{R}$. Then $A_f \neq 0$, hence $A_f$ has a maximal ideal, whose contraction in $A$ is a prime ideal $\mathfrak{p}$ such that $f \notin \mathfrak{p}$ and which is maximal with respect to this property. Then $\mathfrak{p}$ is not maximal and is not equal to the intersection of the prime ideals strictly containing $\mathfrak{p}$.

    A ring $A$ with the three equivalent properties above is called a \emph{Jacobson ring}\index{Jacobson ring}.

    \item \label{exer24-chapter5}Let $A$ be a Jacobson ring (Exercise \ref{exer23-chapter5}) and $B$ an $A$-algebra. Show that if $B$ is either (i) integral over $A$ or (ii) finitely generated as an $A$-algebra, then $B$ is Jacobson.
    
    \textit{Hint:} Use Exercise \ref{exer22-chapter5} for (ii).

    \quad\, In particular, every finitely generated ring, and every finitely generated algebra over a field, is a Jacobson ring.

    \item \label{exer25-chapter5}Let $A$ be a ring. Show that the following are equivalent:
    \begin{enumerate}[i)]
        \item $A$ is a Jacobson ring;
        \item Every finitely generated $A$-algebra $B$ which is a field is finite over $A$.
    \end{enumerate}

    \textit{Hint:} i) $\implies$ ii). Reduce to the case where $A$ is a subring of $B$, and use Exercise \ref{exer21-chapter5}. If $s \in A$ is as in Exercise \ref{exer21-chapter5}, then there exists a maximal ideal $\mathfrak{m}$ of $A$ not containing $s$, and the homomorphism $A \to A/\mathfrak{m} = k$ extends to a homomorphism $g$ of $B$ into the algebraic closure of $k$. Since $B$ is a field, $g$ is injective, and $g(B)$ is algebraic over $k$, hence finite algebraic over $k$.

    ii) $\implies$ i). Use criterion iii) of Exercise \ref{exer23-chapter5}. Let $\mathfrak{p}$ be a prime ideal of $A$ which is not maximal, and let $B = A/\mathfrak{p}$. Let $f$ be a non-zero element of $B$. Then $B_f$ is a finitely generated $A$-algebra. If it is a field it is finite over $B$, hence integral over $B$ and therefore $B$ is a field by \eqref{prop:5.7}. Hence $B_f$ is not a field and therefore has a non-zero prime ideal, whose contraction in $B$ is a non-zero ideal $\mathfrak{p}'$ such that $f \notin \mathfrak{p}'$.

    \item \label{exer26-chapter5}Let $X$ be a topological space. A subset of $X$ is \emph{locally closed}\index{locally closed} if it is the intersection of an open set and a closed set, or equivalently if it is open in its closure.
    
    \quad\, The following conditions on a subset $X_0$ of $X$ are equivalent:
    \begin{enumerate}
        \item[(1)] Every non-empty locally closed subset of $X$ meets $X_0$;
        \item[(2)] For every closed set $E$ in $X$ we have $\overline{E \cap X_0} = \overline{E}$;
        \item[(3)] The mapping $U \mapsto U \cap X_0$ of the collection of open sets of $X$ onto the collection of open sets of $X_0$ is bijective.
    \end{enumerate}
    \quad\, A subset $X_0$ satisfying these conditions is said to be \emph{very dense}\index{very dense} in $X$.

    \quad\, If $A$ is a ring, show that the following are equivalent:
    \begin{enumerate}[i)]
        \item $A$ is a Jacobson ring;
        \item The set of maximal ideals of $A$ is very dense in $\Spec(A)$;
        \item Every locally closed subset of $\Spec(A)$ consisting of a single point is closed.
    \end{enumerate}
    \textit{Hint:} (ii) and (iii) are geometrical formulations of conditions ii) and iii) of \eqref{exer23-chapter5}.

    \emph{Valuation rings and valuations.}

    \item \label{exer27-chapter5} Let $A, B$ be two local rings. $B$ is said to \emph{dominate}\index{domination} $A$ if $A$ is a subring of $B$ and the maximal ideal $\mathfrak{m}$ of $A$ is contained in the maximal ideal $\mathfrak{n}$ of $B$ (or, equivalently, if $\mathfrak{m} = \mathfrak{n} \cap A$). Let $K$ be a field and let $\Sigma$ be the set of all local subrings of $K$. If $\Sigma$ is ordered by the relation of domination, show that $\Sigma$ has maximal elements and that $A \in \Sigma$ is maximal if and only if $A$ is a valuation ring of $K$.
    
    \textit{Hint:} Use \eqref{thm:5.21}.

    \item \label{exer28-chapter5}Let $A$ be an integral domain, $K$ its field of fractions. Show that the following are equivalent:
    \begin{enumerate}
        \item[(1)] $A$ is a valuation ring of $K$;
        \item[(2)] If $\mathfrak{a}, \mathfrak{b}$ are any two ideals of $A$, then either $\mathfrak{a} \subseteq \mathfrak{b}$ or $\mathfrak{b} \subseteq \mathfrak{a}$.
    \end{enumerate}
    \quad\, Deduce that if $A$ is a valuation ring and $\mathfrak{p}$ is a prime ideal of $A$, then $A_{\mathfrak{p}}$ and $A/\mathfrak{p}$ are valuation rings of their fields of fractions.

    \item \label{exer29-chapter5}Let $A$ be a valuation ring of a field $K$. Show that every subring of $K$ which contains $A$ is a local ring of $A$.

    \item \label{exer30-chapter5}Let $A$ be a valuation ring of a field $K$. The group $U$ of units of $A$ is a subgroup of the multiplicative group $K^*$ of $K$.
    
    \quad\, Let $\Gamma = K^*/U$. If $\xi, \eta \in \Gamma$ are represented by $x, y \in K$, define $\xi \geq \eta$ to mean $xy^{-1} \in A$. Show that this defines a total ordering on $\Gamma$ which is compatible with the group structure (i.e., $\xi \geq \eta \implies \xi\omega \geq \eta\omega$ for all $\omega \in \Gamma$). In other words, $\Gamma$ is a totally ordered abelian group. It is called the \emph{value group}\index{value group} of $A$.
    
    Let $v \colon K^* \to \Gamma$ be the canonical homomorphism. Show that $v(x + y) \geq \min(v(x), v(y))$ for all $x, y \in K^*$.

    \item \label{exer31-chapter5}Conversely, let $\Gamma$ be a totally ordered abelian group (written additively), and let $K$ be a field. A \emph{valuation}\index{valuation} of $K$ with values in $\Gamma$ is a mapping $v \colon K^* \to \Gamma$ such that
    \begin{enumerate}
        \item[(1)] $v(xy) = v(x) + v(y)$,
        \item[(2)] $v(x + y) \geq \min(v(x), v(y))$,
    \end{enumerate}
    for all $x, y \in K^*$. Show that the set of elements $x \in K^*$ such that $v(x) \geq 0$ is a valuation ring of $K$. This ring is called the valuation ring of $v$, and the subgroup $v(K^*)$ of $\Gamma$ is the value group of $v$.\\
    Thus the concepts of valuation ring and valuation are essentially equivalent.

    \item \label{exer32-chapter5}Let $\Gamma$ be a totally ordered abelian group. A subgroup $\Delta$ of $\Gamma$ is \emph{isolated}\index{isolated subgroup} in $\Gamma$ if, whenever $0 \leq \beta \leq \alpha$ and $\alpha \in \Delta$, we have $\beta \in \Delta$. Let $A$ be a valuation ring of a field $K$, with value group $\Gamma$ (Exercise 31). If $\mathfrak{p}$ is a prime ideal of $A$, show that $v(A \setminus \mathfrak{p})$ is the set of elements $\geq 0$ in an isolated subgroup $\Delta$ of $\Gamma$, and that the mapping so defined of $\Spec(A)$ into the set of isolated subgroups of $\Gamma$ is bijective.
    
    \quad\, If $\mathfrak{p}$ is a prime ideal of $A$, what are the value groups of the valuation rings $A/\mathfrak{p}$, $A_{\mathfrak{p}}$?

    \item \label{exer33-chapter5}Let $\Gamma$ be a totally ordered abelian group. We shall show how to construct a field $K$ and a valuation $v$ of $K$ with $\Gamma$ as value group. Let $k$ be any field and let $A = k[\Gamma]$ be the group algebra of $\Gamma$ over $k$. By definition, $A$ is freely generated as a $k$-vector space by elements $X_{\alpha}$ ($\alpha \in \Gamma$) such that $X_{\alpha}X_{\beta} = X_{\alpha+\beta}$. Show that $A$ is an integral domain.
    
    \quad\, If $u = \lambda_1 X_{\alpha_1} + \cdots + \lambda_n X_{\alpha_n}$ is any non-zero element of $A$, where the $\lambda_i$ are all $\neq 0$ and $\alpha_1 < \cdots < \alpha_n$, define $v_0(u)$ to be $\alpha_1$. Show that the mapping $v_0 \colon A \setminus \{0\} \to \Gamma$ satisfies conditions (1) and (2) of Exercise \ref{exer31-chapter5}.

    \quad\, Let $K$ be the field of fractions of $A$. Show that $v_0$ can be uniquely extended to a valuation $v$ of $K$, and that the value group of $v$ is precisely $\Gamma$.

    \item \label{exer34-chapter5}Let $A$ be a valuation ring and $K$ its field of fractions. Let $f \colon A \to B$ be a ring homomorphism such that $f^* \colon \Spec(B) \to \Spec(A)$ is a closed mapping. Then if $g \colon B \to K$ is any $A$-algebra homomorphism (i.e., if $g \circ f$ is the embedding of $A$ in $K$) we have $g(B) = A$.\\
    \textit{Hint:} Let $C = g(B)$; obviously $C \supseteq A$. Let $\mathfrak{n}$ be a maximal ideal of $C$. Since $f^*$ is closed, $\mathfrak{m} = \mathfrak{n} \cap A$ is the maximal ideal of $A$, whence $A_{\mathfrak{m}} = A$. Also the local ring $C_{\mathfrak{n}}$ dominates $A_{\mathfrak{m}}$. Hence by Exercise 27 we have $C_{\mathfrak{n}} = A$ and therefore $C \subseteq A$.

    \item \label{exer35-chapter5}From Exercises \ref{exer1-chapter5} and \ref{exer3-chapter5} it follows that, if $f \colon A \to B$ is integral and $C$ is any $A$-algebra, then the mapping $(f \otimes 1)^* \colon \Spec(B \otimes_A C) \to \Spec(C)$ is a closed map.
    
    \quad\, Conversely, suppose that $f \colon A \to B$ has this property and that $B$ is an integral domain. Then $f$ is integral.

    \quad\, \textit{Hint:} Replacing $A$ by its image in $B$, reduce to the case where $A \subseteq B$ and $f$ is the injection. Let $K$ be the field of fractions of $B$ and let $A'$ be a valuation ring of $K$ containing $A$. By \eqref{cor:5.22} it is enough to show that $A'$ contains $B$. By hypothesis $\Spec(B \otimes_A A') \to \Spec(A')$ is a closed map. Apply the result of Exercise \ref{exer34-chapter5} to the homomorphism $B \otimes_A A' \to K$ defined by $b \otimes a' \mapsto ba'$. It follows that $ba' \in A'$ for all $b \in B$ and all $a' \in A'$; taking $a' = 1$, we have what we want.

    \quad\, Show that the result just proved remains valid if $B$ is a ring with only finitely many minimal prime ideals (e.g., if $B$ is Noetherian).

    \quad\, \textit{Hint:} Let $\mathfrak{p}_i$ be the minimal prime ideals. Then each composite homomorphism $A \to B \to B/\mathfrak{p}_i$ is integral, hence $A \to \prod (B/\mathfrak{p}_i)$ is integral, hence $A \to B/\mathfrak{N}$ is integral (where $\mathfrak{N}$ is the nilradical of $B$), hence finally $A \to B$ is integral.
\end{enumerate}

